\documentclass[11pt,a4paper]{article}

% ── Packages ────────────────────────────────────────────────────────────────
\usepackage[margin=1in]{geometry}
\usepackage{amsmath,amssymb,amsfonts}
\usepackage{booktabs}
\usepackage{array}
\usepackage{xcolor}
\usepackage{hyperref}
\usepackage{caption}
\usepackage{float}

% ── Math operators ──────────────────────────────────────────────────────────
\DeclareMathOperator{\meanpool}{mean\_pool}
\DeclareMathOperator{\stack}{stack}
\DeclareMathOperator{\unfold}{unfold}
\DeclareMathOperator{\diag}{diag}

% ── Title ───────────────────────────────────────────────────────────────────
\title{Population HOSVD Hallucination Detection Pipeline\\
       \large Mathematical Formulation \& Dimensional Audit}
\author{Graduate Research — Mechanistic Interpretability}
\date{}

\begin{document}
\maketitle

% ══════════════════════════════════════════════════════════════════════════════
\section{Anatomical Constants}
% ══════════════════════════════════════════════════════════════════════════════

\begin{table}[H]
\centering
\begin{tabular}{ccl}
\toprule
\textbf{Symbol} & \textbf{Value} & \textbf{Meaning} \\
\midrule
$L$     & 32    & Transformer layers \\
$D$     & 4096  & Hidden dimension (residual-stream width) \\
$N$     & 3979  & Training population size \\
$R_L$   & 5     & Target rank for layer (Mode-1) \\
$R_D$   & 64    & Target rank for hidden (Mode-2) \\
$T_i$   & $\in [4,16]$ & Variable token count per prompt \\
\bottomrule
\end{tabular}
\caption{Global constants for the Phase~I LLaMA-2~7B experiment.}
\end{table}

% ══════════════════════════════════════════════════════════════════════════════
\section{Pipeline Stages}
% ══════════════════════════════════════════════════════════════════════════════

Each stage is presented as a triple:
\begin{enumerate}
  \item \textbf{Operation} — the linear-algebraic step.
  \item \textbf{Dimensional audit} — exact tensor shapes before and after.
  \item \textbf{Running example} — a concrete numeric trace for a single sample.
\end{enumerate}

% ──────────────────────────────────────────────────────────────────────────────
\subsection{Stage 1 — Raw Residual-Stream Harvest}
% ──────────────────────────────────────────────────────────────────────────────

\textbf{Concept.}
A forward pass of a transformer on prompt $i$ yields hidden states at every layer.
We collect these into a 3-mode tensor.

\[
\mathcal{X}^{(i)} \in \mathbb{R}^{L \times T_i \times D}
\]

\noindent
\textbf{Dimensions.}

\begin{table}[H]
\centering
\begin{tabular}{>{\ttfamily}l>{\ttfamily}l}
\toprule
\multicolumn{1}{c}{\textbf{Tensor}} & \multicolumn{1}{c}{\textbf{Shape}} \\
\midrule
raw\_activation\_list[i]  & (32, $T_i$, 4096) \\
raw\_activation\_list     & list of $N$ such tensors \\
\bottomrule
\end{tabular}
\caption{Stage~1 — raw activations per sample.}
\end{table}

\noindent
\textbf{Running example \#1.}
Sample~0 has $T_0 = 14$ tokens (a medium-length prompt).
\[
\mathcal{X}^{(0)} \in \mathbb{R}^{32 \times 14 \times 4096}
\]
For each of the 32 layers we have a $(14 \times 4096)$ matrix: 14 token
representations, each a 4096-dimensional vector.

% ──────────────────────────────────────────────────────────────────────────────
\subsection{Stage 2 — Token Neutralization (Mean-Pooling)}
% ──────────────────────────────────────────────────────────────────────────────

\textbf{Operation.}
The token axis (Mode-1) is collapsed by taking the arithmetic mean,
discarding positional information but preserving the average activation
geometry per layer.

\[
\overline{\mathcal{X}}^{(i)}
   = \frac{1}{T_i} \sum_{t=1}^{T_i} \mathcal{X}^{(i)}_{:,t,:}
   \;\in\; \mathbb{R}^{L \times D}
\]

All $N$ pooled matrices are stacked along a new population axis (Mode-0):

\[
\mathbf{X}_{\text{pop}} = \bigl[\,\overline{\mathcal{X}}^{(1)},\;
                              \overline{\mathcal{X}}^{(2)},\; \dots,\;
                              \overline{\mathcal{X}}^{(N)}\,\bigr]
                        \;\in\; \mathbb{R}^{N \times L \times D}
\]

\noindent
\textbf{Dimensions.}

\begin{table}[H]
\centering
\begin{tabular}{>{\ttfamily}l>{\ttfamily}l}
\toprule
\multicolumn{1}{c}{\textbf{Tensor}} & \multicolumn{1}{c}{\textbf{Shape}} \\
\midrule
$\overline{\mathcal{X}}^{(i)}$ (per sample) & (32, 4096) \\
$\mathbf{X}_{\text{pop}}$ (stacked) & (3979, 32, 4096) \\
\bottomrule
\end{tabular}
\caption{Stage~2 — mean-pooled and stacked population tensor.}
\end{table}

\noindent
\textbf{Running example \#2.}
Sample~0's 14 tokens are averaged into a single $(32 \times 4096)$ matrix.
After stacking all 3979 samples we obtain
\[
\mathbf{X}_{\text{pop}} \in \mathbb{R}^{3979 \times 32 \times 4096}.
\]

% ──────────────────────────────────────────────────────────────────────────────
\subsection{Stage 3 — Gram Matrix Trick (Factor Matrices $U_L$ and $U_D$)}
% ──────────────────────────────────────────────────────────────────────────────

\textbf{Problem.}
A direct SVD on the unfolded population tensor would require factorising a
$32 \times (3979\cdot 4096) = 32 \times 16\,302\,784$ matrix — wasteful.
Instead we eigen-decompose the \emph{Gram matrix} $X X^{\mathsf{T}}$,
which is only $32 \times 32$.

\medskip
\noindent
\textbf{Step~A — Layer factor matrix $U_L$.}

Unfold $\mathbf{X}_{\text{pop}}$ along Mode~1 (layers):
\[
\mathbf{X}_{(L)} = \unfold_{(1)}\!\bigl(\mathbf{X}_{\text{pop}}\bigr)
                 \in \mathbb{R}^{L \times (N \cdot D)}
                 = \mathbb{R}^{32 \times 16\,302\,784}
\]

Form the symmetric positive-semidefinite Gram matrix and eigen-decompose:
\[
\mathbf{A}_L = \mathbf{X}_{(L)} \, \mathbf{X}_{(L)}^{\mathsf{T}}
             \in \mathbb{R}^{L \times L}
             = \mathbb{R}^{32 \times 32}
\]
\[
\mathbf{A}_L = \mathbf{V}_L \, \Lambda_L \, \mathbf{V}_L^{\mathsf{T}},
\qquad
\Lambda_L = \diag(\lambda_1,\dots,\lambda_{32}),\quad
\lambda_1 \ge \lambda_2 \ge \dots \ge \lambda_{32}
\]

Retain the top $R_L = 5$ eigenvectors (columns of $\mathbf{V}_L$ corresponding
to the 5 largest eigenvalues):
\[
U_L = \mathbf{V}_L[:,\; 1\!:\!R_L] \;\in\; \mathbb{R}^{L \times R_L}
     = \mathbb{R}^{32 \times 5}
\]

\medskip
\noindent
\textbf{Step~B — Hidden-dimension factor matrix $U_D$.}

Unfold along Mode~2 (hidden dimension):
\[
\mathbf{X}_{(D)} = \unfold_{(2)}\!\bigl(\mathbf{X}_{\text{pop}}\bigr)
                 \in \mathbb{R}^{D \times (N \cdot L)}
                 = \mathbb{R}^{4096 \times 127\,328}
\]

Gram matrix:
\[
\mathbf{A}_D = \mathbf{X}_{(D)} \, \mathbf{X}_{(D)}^{\mathsf{T}}
             \in \mathbb{R}^{D \times D}
             = \mathbb{R}^{4096 \times 4096}
\]

Eigen-decompose and keep top $R_D = 64$ eigenvectors:
\[
U_D = \mathbf{V}_D[:,\; 1\!:\!R_D] \;\in\; \mathbb{R}^{D \times R_D}
     = \mathbb{R}^{4096 \times 64}
\]

\noindent
\textbf{Dimensions.}

\begin{table}[H]
\centering
\begin{tabular}{>{\ttfamily}l>{\ttfamily}l}
\toprule
\multicolumn{1}{c}{\textbf{Tensor}} & \multicolumn{1}{c}{\textbf{Shape}} \\
\midrule
$\mathbf{X}_{(L)}$ (mode-1 unfold)  & (32, $3979\cdot 4096$) = (32, 16\,302\,784) \\
$\mathbf{A}_L$ (Gram)               & (32, 32) \\
$U_L$ (top $R_L$ eigenvectors)      & (32, 5) \\[4pt]
$\mathbf{X}_{(D)}$ (mode-2 unfold)  & (4096, $3979\cdot 32$) = (4096, 127\,328) \\
$\mathbf{A}_D$ (Gram)               & (4096, 4096) \\
$U_D$ (top $R_D$ eigenvectors)      & (4096, 64) \\
\bottomrule
\end{tabular}
\caption{Stage~3 — Gram matrix trick and factor matrices.}
\end{table}

\noindent
\textbf{Running example \#3.}
The Gram matrix $\mathbf{A}_L$ is only $32 \times 32$ — eigen-decomposition
completes in $<1$~ms.  $\mathbf{A}_D$ is $4096 \times 4096$ (~67~MB in float32),
still tractable on CPU.  The output bases $U_L$ and $U_D$ are orthonormal:
$U_L^{\mathsf{T}} U_L = I_{5}$ and $U_D^{\mathsf{T}} U_D = I_{64}$.

% ──────────────────────────────────────────────────────────────────────────────
\subsection{Stage 4 — Multilinear Core Compression}
% ──────────────────────────────────────────────────────────────────────────────

\textbf{Operation.}
Each pooled sample matrix $\overline{\mathcal{X}}^{(i)} \in \mathbb{R}^{L \times D}$
is projected onto the factor bases via a two-sided contraction:

\[
\mathcal{G}^{(i)} = U_L^{\mathsf{T}} \;
                    \overline{\mathcal{X}}^{(i)} \;
                    U_D
                  \;\in\; \mathbb{R}^{R_L \times R_D}
                  = \mathbb{R}^{5 \times 64}
\]

Dimensionally:
\[
\underbrace{(5 \times 32)}_{U_L^{\mathsf{T}}}
\cdot
\underbrace{(32 \times 4096)}_{\overline{\mathcal{X}}^{(i)}}
\cdot
\underbrace{(4096 \times 64)}_{U_D}
=
(5 \times 64)
\]

Stacking all $N$ core tensors:
\[
\mathbf{G}_{\text{pop}} = \bigl[\,\mathcal{G}^{(1)},\; \dots,\;
                            \mathcal{G}^{(N)}\,\bigr]
                        \;\in\; \mathbb{R}^{N \times R_L \times R_D}
                        = \mathbb{R}^{3979 \times 5 \times 64}
\]

\noindent
\textbf{Dimensions.}

\begin{table}[H]
\centering
\begin{tabular}{>{\ttfamily}l>{\ttfamily}l}
\toprule
\multicolumn{1}{c}{\textbf{Tensor}} & \multicolumn{1}{c}{\textbf{Shape}} \\
\midrule
$\mathcal{G}^{(i)}$ (per-sample core) & (5, 64) \\
$\mathbf{G}_{\text{pop}}$ (stacked)   & (3979, 5, 64) \\
\bottomrule
\end{tabular}
\caption{Stage~4 — core tensor compression.}
\end{table}

\noindent
\textbf{Running example \#4.}
Sample~0's original $(32 \times 14 \times 4096)$ tensor (14 tokens $\times$ 4096
features $\times$ 32 layers) is compressed to a $(5 \times 64) = 320$-entry core.
This is a compression ratio of:
\[
\frac{32 \cdot 14 \cdot 4096}{5 \cdot 64}
= \frac{1\,835\,008}{320} \approx 5734\times
\]

% ──────────────────────────────────────────────────────────────────────────────
\subsection{Stage 5 — Centroid Computation}
% ──────────────────────────────────────────────────────────────────────────────

\textbf{Operation.}
Split $\mathbf{G}_{\text{pop}}$ by ground-truth binary label $y_i \in \{0,1\}$
($0 =$ truthful, $1 =$ hallucinated).  Compute the arithmetic mean in core space:

\[
\mathcal{C}_{\text{truth}} = \frac{1}{N_0}
      \sum_{i: y_i=0} \mathcal{G}^{(i)}
      \;\in\; \mathbb{R}^{R_L \times R_D}
      = \mathbb{R}^{5 \times 64}
\]
\[
\mathcal{C}_{\text{hall}} = \frac{1}{N_1}
      \sum_{i: y_i=1} \mathcal{G}^{(i)}
      \;\in\; \mathbb{R}^{R_L \times R_D}
      = \mathbb{R}^{5 \times 64}
\]

where $N_0 = 2341$ (truthful) and $N_1 = 1638$ (hallucinated).

\noindent
\textbf{Dimensions.}

\begin{table}[H]
\centering
\begin{tabular}{>{\ttfamily}l>{\ttfamily}l}
\toprule
\multicolumn{1}{c}{\textbf{Tensor}} & \multicolumn{1}{c}{\textbf{Shape}} \\
\midrule
$\mathcal{C}_{\text{truth}}$ (truth centroid)   & (5, 64) \\
$\mathcal{C}_{\text{hall}}$ (halluc. centroid)  & (5, 64) \\
\bottomrule
\end{tabular}
\caption{Stage~5 — class centroids in compressed core space.}
\end{table}

\noindent
\textbf{Running example \#5.}
The two centroids are separated by a Frobenius distance of
$\|\mathcal{C}_{\text{truth}} - \mathcal{C}_{\text{hall}}\|_F \approx 19.5$.
The dominant Mode~0 alone accounts for $17.9$ ($\approx 92\%$) of this distance.

% ──────────────────────────────────────────────────────────────────────────────
\subsection{Stage 6 — Test-Time Detective}
% ──────────────────────────────────────────────────────────────────────────────

\textbf{Operation.}
A held-out test activation $\mathcal{X}^{(\text{test})} \in \mathbb{R}^{L \times T_{\text{test}} \times D}$
is passed through the identical pipeline using the \emph{training-derived}
factor matrices $U_L, U_D$:

\begin{align}
\overline{\mathcal{X}}^{(\text{test})}
   &= \meanpool\!\bigl(\mathcal{X}^{(\text{test})}\bigr)
      \;\in\; \mathbb{R}^{L \times D}
      = \mathbb{R}^{32 \times 4096} \\[4pt]
\mathcal{G}^{(\text{test})}
   &= U_L^{\mathsf{T}} \; \overline{\mathcal{X}}^{(\text{test})} \; U_D
      \;\in\; \mathbb{R}^{R_L \times R_D}
      = \mathbb{R}^{5 \times 64}
\end{align}

Four scalar features are extracted:

\begin{align}
d_C &= \bigl\|\mathcal{G}^{(\text{test})} - \mathcal{C}_{\text{truth}}\bigr\|_F
      & &\text{(distance to truth centroid)} \\[2pt]
d_H &= \bigl\|\mathcal{G}^{(\text{test})} - \mathcal{C}_{\text{hall}}\bigr\|_F
      & &\text{(distance to hallucination centroid)} \\[2pt]
\Delta &= d_C - d_H
      & &\text{(signed margin: $-$ $\rightarrow$ truth, $+$ $\rightarrow$ halluc.)} \\[2pt]
\eta   &= \bigl\|\mathcal{G}^{(\text{test})}\bigr\|_F
      & &\text{(core activation magnitude)}
\end{align}

\noindent
\textbf{Dimensions.}

\begin{table}[H]
\centering
\begin{tabular}{>{\ttfamily}l>{\ttfamily}l}
\toprule
\multicolumn{1}{c}{\textbf{Tensor}} & \multicolumn{1}{c}{\textbf{Shape}} \\
\midrule
$\mathcal{X}^{(\text{test})}$ (raw)           & (32, $T_{\text{test}}$, 4096) \\
$\overline{\mathcal{X}}^{(\text{test})}$      & (32, 4096) \\
$\mathcal{G}^{(\text{test})}$                 & (5, 64) \\
$d_C,\; d_H,\; \Delta,\; \eta$ (features)     & scalar \\
\bottomrule
\end{tabular}
\caption{Stage~6 — test-time feature extraction.}
\end{table}

\noindent
\textbf{Running example \#6 (held-out sample).}
\begin{align*}
d_C    &= 98.39 \quad\text{(far from truth)} \\
d_H    &= 83.43 \quad\text{(closer to hallucination)} \\
\Delta &= +14.96 \;\rightarrow\; \text{predicted hallucination} \\
\eta   &= 451.53
\end{align*}

% ──────────────────────────────────────────────────────────────────────────────
\subsection{Stage 7 — Supervised Classification}
% ──────────────────────────────────────────────────────────────────────────────

\textbf{Operation.}
The core tensors are flattened into a 2D feature matrix and fed to
standard classifiers:

\[
\mathbf{X}_{\text{flat}} = \operatorname{reshape}\!\bigl(
      \mathbf{G}_{\text{pop}},\; (N,\; R_L \cdot R_D)\bigr)
      \;\in\; \mathbb{R}^{N \times (R_L \cdot R_D)}
      = \mathbb{R}^{3979 \times 320}
\]

Each row of $\mathbf{X}_{\text{flat}}$ is the 320-dimensional feature vector
for one sample.  An 80/20 stratified train/test split is applied, and two
classifiers are trained:

\begin{itemize}
  \item \textbf{Logistic Regression} — linear hyperplane in core space.
  \item \textbf{Random Forest} (200 trees, balanced class weights) — non-linear
        decision boundary.
\end{itemize}

\noindent
\textbf{Dimensions.}

\begin{table}[H]
\centering
\begin{tabular}{>{\ttfamily}l>{\ttfamily}l}
\toprule
\multicolumn{1}{c}{\textbf{Tensor}} & \multicolumn{1}{c}{\textbf{Shape}} \\
\midrule
$\mathbf{G}_{\text{pop}}$ (3D cores)  & (3979, 5, 64) \\
$\mathbf{X}_{\text{flat}}$ (2D features) & (3979, 320) \\
$X_{\text{train}}$ / $X_{\text{test}}$   & (3183, 320) / (796, 320) \\
\bottomrule
\end{tabular}
\caption{Stage~7 — flattening and classification.}
\end{table}

\noindent
\textbf{Results (Phase~I, LLaMA-2~7B, 80/20 split).}

\begin{table}[H]
\centering
\begin{tabular}{lc}
\toprule
\textbf{Classifier} & \textbf{AUROC} \\
\midrule
Logistic Regression  & 0.825 \\
Random Forest        & 0.958 \\
\bottomrule
\end{tabular}
\caption{Classification performance on Phase~I data (no data leakage).}
\end{table}


% ══════════════════════════════════════════════════════════════════════════════
\section{Cross-Model Generalization: Phase II (Qwen~2.5--3B)}
% ══════════════════════════════════════════════════════════════════════════════

To assess whether the HOSVD subspace projection generalises beyond a single
model, we replicate the full pipeline on a second, smaller LLM with
independently generated labels.

\subsection{Experimental Design}

\begin{table}[H]
\centering
\begin{tabular}{lcc}
\toprule
\textbf{Property} & \textbf{Phase I} & \textbf{Phase II} \\
\midrule
Model                & LLaMA-2~7B        & Qwen~2.5--3B-Instruct \\
Precision            & (pre-extracted)   & bfloat16 native \\
Layers ($L$)         & 32                & 36 \\
Hidden dim ($D$)     & 4096              & 2048 \\
Samples ($N$)        & 3979              & 817 \\
Prompt source        & Custom corpus     & TruthfulQA (validation) \\
Label source         & Pre-annotated     & BLEURT-20 $\lor$ ROUGE-L \\
Label threshold      & Binary ground-truth & $\text{BLEURT}\ge 0.5 \lor \text{ROUGE-L}\ge 0.7$ \\
Hallucination rate   & 41.2\%            & 81.4\% \\
Compression target   & $(5, 64)$         & $(5, 64)$ \\
\bottomrule
\end{tabular}
\caption{Phase~I vs.\ Phase~II — experimental configuration.}
\end{table}

\noindent
\textbf{Key differences.}
Phase~II uses a \emph{different} model architecture, \emph{different} prompt
distribution (TruthfulQA is adversarially designed), and \emph{automatic}
label generation via BLEURT/ROUGE-L rather than human annotation.  Any
consistent signal across both phases constitutes evidence that the subspace
method captures a model-agnostic hallucination signature.

\subsection{Core Distribution Comparison}

\begin{table}[H]
\centering
\begin{tabular}{lcccc}
\toprule
\textbf{Metric} & \multicolumn{2}{c}{\textbf{Phase I (LLaMA-2)}} & \multicolumn{2}{c}{\textbf{Phase II (Qwen)}} \\
\cmidrule(lr){2-3} \cmidrule(lr){4-5}
& Truthful & Halluc. & Truthful & Halluc. \\
\midrule
Count                          & 2341 (58.8\%) & 1638 (41.2\%) & 152 (18.6\%) & 665 (81.4\%) \\
Frobenius norm  (mean $\pm$ std) & $377 \pm 216$  & $393 \pm 236$  & $676 \pm 57$  & $660 \pm 72$ \\
Centroid norm                  & 364.4 & 381.3 & 648 & 628 \\
\addlinespace
\textbf{Layer-mode energy}     &        &        &        & \\
\quad Mode 0 (dominant)        & 20\,237 & 17\,929 & 8\,000 & 16\,384 \\
\quad Mode 1                   & 5\,203  & 4\,306  & 1\,648 & 3\,440 \\
\quad Mode 2                   & 2\,124  & 1\,799  & 1\,320 & 2\,672 \\
\quad Mode 3                   & 874     & 740     & 820   & 1\,744 \\
\quad Mode 4                   & 502     & 424     & 548   & 1\,144 \\
\addlinespace
\textbf{Cross-group separation}&        &        &        & \\
\quad Centroid distance        & \multicolumn{2}{c}{19.49} & \multicolumn{2}{c}{50.25} \\
\quad Mode 0 diff              & \multicolumn{2}{c}{17.94 (92\%)} & \multicolumn{2}{c}{44.50 (89\%)} \\
\quad Mode 1 diff              & \multicolumn{2}{c}{6.81}  & \multicolumn{2}{c}{14.13} \\
\quad Mode 2 diff              & \multicolumn{2}{c}{3.17}  & \multicolumn{2}{c}{13.06} \\
\quad Mode 3 diff              & \multicolumn{2}{c}{1.10}  & \multicolumn{2}{c}{11.44} \\
\quad Mode 4 diff              & \multicolumn{2}{c}{0.62}  & \multicolumn{2}{c}{6.22} \\
\bottomrule
\end{tabular}
\caption{Core-space distribution statistics — Phase~I vs.\ Phase~II.}
\end{table}

\noindent
\textbf{Observations.}
\begin{enumerate}
  \item \textbf{Centroid separation persists across models.}
        The inter-class centroid distance is \emph{larger} in Phase~II (50.25
        vs.\ 19.49), despite the smaller dataset.  In both phases, the
        dominant Mode~0 carries $\approx 90\%$ of the separation — the
        highest-variance layer-interaction axis is also the most
        discriminative.

  \item \textbf{Hallucination cores are hotter in the dominant mode.}
        In Phase~II, hallucinated samples have \emph{double} the Mode-0 energy
        of truthful samples (16\,384 vs.\ 8\,000).  This asymmetry is more
        pronounced than in Phase~I, likely because TruthfulQA's adversarial
        prompts force the model into high-activation failure modes.

  \item \textbf{Label imbalance is a confound, not a fatal flaw.}
        Phase~II is heavily imbalanced (81.4\% hallucinated).  AUROC is
        robust to class imbalance, and we use stratified splits and balanced
        class weights throughout.

  \item \textbf{Automatic labels are noisier than human labels.}
        BLEURT/ROUGE-L thresholds introduce label noise, which reduces
        achievable AUROC.  The Phase~I $\to$ Phase~II drop (0.96 $\to$ 0.74)
        is partially attributable to label quality, not just model/size
        differences.
\end{enumerate}

\subsection{Classification Results}

\begin{table}[H]
\centering
\begin{tabular}{lcc}
\toprule
\textbf{Classifier} & \textbf{Phase I (LLaMA-2)} & \textbf{Phase II (Qwen)} \\
\midrule
Logistic Regression  & 0.825 & 0.721 \\
Random Forest        & 0.958 & 0.740 \\
\midrule
Train / test split   & 3183 / 796 & 653 / 164 \\
Feature dimension    & 320        & 320 \\
\bottomrule
\end{tabular}
\caption{No-leakage AUROC — Phase~I vs.\ Phase~II.  Factor matrices
         $U_L$, $U_D$ are computed \emph{exclusively} from the training
         partition in both cases.}
\end{table}

\noindent
\textbf{Interpretation.}
The HOSVD subspace projection generalises across model architectures and
label sources.  A Random Forest operating on 320-dimensional compressed core
vectors achieves $>0.74$ AUROC on both datasets, well above the chance
baseline of 0.50.  The Phase~I $\to$ Phase~II gap (0.96 $\to$ 0.74) is
consistent with the compounding effects of:
\begin{itemize}
  \item Smaller training set (3979 $\to$ 817; a $4.9\times$ reduction);
  \item Smaller model hidden dimension (4096 $\to$ 2048; less expressive
        residual stream);
  \item Noisier automatic labels (BLEURT/ROUGE-L vs.\ human annotation);
  \item Stronger class imbalance requiring heavier regularisation.
\end{itemize}

Critically, the \emph{method} does not degrade to chance — the core tensors
encode a model-agnostic hallucination signature that a non-linear classifier
can recover even under adverse conditions.


% ══════════════════════════════════════════════════════════════════════════════
\section{Complete Dimensional Summary}
% ══════════════════════════════════════════════════════════════════════════════

\begin{table}[H]
\centering
\begin{tabular}{c>{\ttfamily}l>{\ttfamily}l>{\ttfamily}l}
\toprule
\textbf{Stage} & \textbf{Operation} & \textbf{Phase I Shape} & \textbf{Phase II Shape} \\
\midrule
1 & Raw harvest (per sample)           & (32, $T_i$, 4096)   & (36, $T_i$, 2048) \\
2 & Mean-pool \& stack                 & (3979, 32, 4096)    & (817, 36, 2048) \\
3a & Gram $\rightarrow$ $U_L$          & (32, 5)             & (36, 5) \\
3b & Gram $\rightarrow$ $U_D$          & (4096, 64)          & (2048, 64) \\
4 & $U_L^{\mathsf{T}} \cdot x \cdot U_D$ (per sample) & (5, 64) & (5, 64) \\
4 & Stack $\rightarrow \mathbf{G}_{\text{pop}}$ & (3979, 5, 64) & (817, 5, 64) \\
5 & Centroid $\mathcal{C}_{\text{truth}}$, $\mathcal{C}_{\text{hall}}$ & (5, 64) each & (5, 64) each \\
6 & Test projection $\mathcal{G}^{(\text{test})}$ & (5, 64) & (5, 64) \\
7 & Flatten $\rightarrow \mathbf{X}_{\text{flat}}$ & (3979, 320) & (817, 320) \\
\bottomrule
\end{tabular}
\caption{Every tensor shape in the pipeline — Phase~I (LLaMA-2~7B) vs.\ 
         Phase~II (Qwen~2.5--3B).  The core tensors converge to identical
         $(5, 64)$ shape regardless of the source model's anatomy.}
\end{table}

% ══════════════════════════════════════════════════════════════════════════════
\section{Key Equations}
% ══════════════════════════════════════════════════════════════════════════════

\noindent
\textbf{Multilinear projection (core compression).}
\begin{equation}
\mathcal{G}^{(i)}_{a,b}
   = \sum_{l=1}^{L} \sum_{d=1}^{D}
     (U_L)_{l,a} \;
     \overline{\mathcal{X}}^{(i)}_{l,d} \;
     (U_D)_{d,b}
   \qquad
   a = 1,\dots,R_L,\;\; b = 1,\dots,R_D
\end{equation}

\noindent
\textbf{Gram matrix trick (avoiding large SVD).}
\begin{equation}
\mathbf{A}_L = \mathbf{X}_{(L)} \mathbf{X}_{(L)}^{\mathsf{T}},\qquad
\mathbf{A}_L \mathbf{v}_k = \lambda_k \mathbf{v}_k,\qquad
U_L = [\mathbf{v}_{L-R_L+1},\dots,\mathbf{v}_{L}]
\end{equation}
(The eigenvectors are flipped after \texttt{eigh} so Mode~0 corresponds to
the largest eigenvalue.)

\noindent
\textbf{Frobenius distance (test-time metric).}
\begin{equation}
d_C = \Bigl(\sum_{a=1}^{R_L}\sum_{b=1}^{R_D}
      \bigl(\mathcal{G}^{(\text{test})}_{a,b}
          - (\mathcal{C}_{\text{truth}})_{a,b}\bigr)^2
      \Bigr)^{1/2}
\end{equation}

\noindent
\textbf{Compression ratio.}
\begin{equation}
\rho = \frac{L \cdot \bar{T} \cdot D}{R_L \cdot R_D}
     \approx \frac{32 \cdot 10 \cdot 4096}{5 \cdot 64}
     = \frac{1\,310\,720}{320} \approx 4096\times
\end{equation}

\noindent
\textbf{Per-mode energy decomposition.}
\begin{equation}
E_k = \Bigl(\sum_{i=1}^{N} \sum_{b=1}^{R_D}
      \bigl(\mathcal{G}^{(i)}_{k,b}\bigr)^2\Bigr)^{1/2},
\qquad k = 0,\dots,R_L-1
\end{equation}

\end{document}
